% Prova 01 remediation --- practice list on GDP measurement (error cluster 1 of Prova/prova-01-error-analysis.md).
% \solucoesfalse -> blank list ; \solucoestrue -> answer key + rubric (copy saved as practice-P1-gdp-gabarito.tex).
% Compile from this folder: pdflatex (2x); check: pdffonts (all Type 1). Every number: practice-P1-gdp-check.py
\documentclass[11pt,a4paper]{article}
% Shared preamble of the practice lists (practice-L*.tex), adapted from lista-exam.tex.
% Each list sets \solucoesfalse / \solucoestrue and \listatema right after \input{practice-preamble}.
\usepackage[utf8]{inputenc}
\usepackage[T1]{fontenc}
\usepackage{lmodern}   % Type 1 vector fonts (replaces PK bitmaps)
\usepackage{cmap}      % ToUnicode CMap -> copyable text and accents
\usepackage[english]{babel}
\usepackage{amsmath,amssymb,amsfonts}
\usepackage[a4paper,margin=2.2cm]{geometry}
\usepackage{enumitem}
\usepackage{fancyhdr}
\usepackage{comment}
% One cut file per job, so builds running side by side cannot swap each other's solutions.
\def\DefaultCutFileName{\def\CommentCutFile{\jobname.cut}}
\usepackage{xcolor}
\usepackage{microtype}
\usepackage{booktabs}
\usepackage{tikz}
\usepackage{pgfplots}
\pgfplotsset{compat=1.16}

\newif\ifsolucoes
\newcommand{\disciplina}{Macroeconomics I --- MPE Insper}
\newcommand{\espacoresol}{4.5cm}   % height of the answer space in the blank list
% The sol environment must be decided after the list sets its toggle.
\AtBeginDocument{\ifsolucoes\includecomment{sol}\else\excludecomment{sol}\fi}

\pagestyle{fancy}\fancyhf{}
\lhead{\footnotesize\disciplina}
\rhead{\footnotesize Practice list\ifsolucoes{} --- Solutions and rubric\fi}
\cfoot{\footnotesize\thepage}
\renewcommand{\headrulewidth}{0.3pt}
\setlength{\headheight}{14pt}
\setlist{topsep=2pt,itemsep=2pt}
\newcounter{questao}

\newcommand{\bloco}[3]{\par\bigskip\noindent
  {\large\bfseries\color{blue!55!black}Block #1 --- #2}\par\noindent\hfill\textit{(#3)}\par
  \noindent\rule{\linewidth}{0.3pt}\smallskip}
\newcommand{\questao}[1]{\par\refstepcounter{questao}\medskip\noindent
  \textbf{Question \thequestao.}~#1\par}
\newenvironment{alternativas}{\begin{enumerate}[label=\Alph*.,leftmargin=2.2em,topsep=1pt]}{\end{enumerate}}
\newcommand{\gabarito}[1]{\ifsolucoes\par\noindent\textbf{Answer:}~#1\fi}
\newcommand{\resolespaco}{\ifsolucoes\else\par\textit{Solution:}\vspace{\espacoresol}\par\fi}
\newcommand{\passo}[2]{\par\smallskip\noindent\textbf{Step --- #1.}~#2}
\newcommand{\intuicao}[1]{\par\smallskip\noindent\textit{Economic intuition.}~#1}
\newcommand{\refbib}[1]{\par\smallskip\noindent{\footnotesize\color{black!60}\textbf{Ref:} #1}}
\newcommand{\rubrica}[1]{\par\smallskip\noindent\textbf{Grading rubric.}\begin{itemize}[leftmargin=1.4em,topsep=1pt]#1\end{itemize}}

\newcommand{\cabecalho}{%
\begin{center}
  {\Large\bfseries \disciplina}\\[2pt]
  {\large Practice list: \listatema}\ifsolucoes\\[2pt]\textit{Solutions and grading rubric}\fi
\end{center}
\vspace{2pt}
{\footnotesize\noindent\textbf{Instructions.} Multiple choice: mark the single correct option (no
penalty for wrong answers). \emph{True/False}: every 2 wrong items cancel 1 right one
(score $=\max\{0,\text{right}-\text{wrong}/2\}\times$ item value); blank items do not count.
Discursive questions: show every step. Parameters are symbols with the stated admissible
ranges; answers are closed forms, conditions, signs and limits.\par}}

\usepackage{needspace}
\newenvironment{notacao}{\par\smallskip\noindent\textbf{Notation.}\begin{itemize}[leftmargin=1.4em,topsep=1pt,itemsep=0pt]\small}{\end{itemize}}
\solucoesfalse                % <<<< \solucoestrue for the ANSWER KEY
\newcommand{\listatema}{GDP measurement --- three approaches that must agree (Prova 01, Question 1)}
\renewcommand{\espacoresol}{6cm}
\setlength{\emergencystretch}{3em}   % lets the long \texttt URL in Kurlat 1.4 break cleanly

\begin{document}
\cabecalho

{\footnotesize\noindent\textbf{How to use this list.} Book exercises are reproduced verbatim, in the book's own notation and
currency, after a source line; parts this list adds are marked. Every numerical item asks for GDP by the \emph{product} (value-added),
\emph{income} and \emph{expenditure} approaches; an answer that gives one approach only is incomplete, and the
three totals must agree. Lay out each answer as Kurlat does: one table with columns Production $|$ Income $|$
Expenditure. Notation: $Y=C+I+G+X-M$ (Kurlat 2020, eq.~1.1.1, p.~15); A = input firm, B = final-goods firm,
N = ad-financed social network, as in Prova~01. Amounts without a unit are in reais.\par}

% =============================================================================
\bloco{1}{Exam level: the three approaches, item by item}{Kurlat \S1.1; Jones ch.~2 \,\textbullet\, 30 points}

\questao{(20 points) \textit{Source: Kurlat, P. (2020), \emph{A Course in Modern Macroeconomics}, Exercise 1.1 ``Accounting'',
p.~27 --- statement reproduced verbatim, all items (a)--(g).}}
\textit{What this trains:} building the Production\,$|$\,Income\,$|$\,Expenditure table for imports, inventories,
government output valued at cost, a retail margin, investment goods and taxes --- the exact skill Prova~01
Q1(b)--(c) asked for.

\noindent How does GDP accounting record the following events? For each of them, describe how they would be computed in GDP
accounts using the income method, the production method and the expenditure method.
\begin{enumerate}[label=(\alph*),leftmargin=2.2em,topsep=2pt,itemsep=1pt]
\item A car manufacturer buys components from Japan for \$1 to be used in production later on and stores them at its warehouse.
\item A car manufacturer buys components from Japan for \$1 and uses half of those components in the production of a car, that it
sells to Andy for \$2. It stores the rest of the components.
\item An army battalion is deployed to the border to repel a threatened Canadian invasion. The soldiers earn wages of \$10,000 and
use ammunition that the government buys for \$5,000. The ammunition is produced using \$2,000 of imported steel and 100 hours of
work, for which the workers were paid \$1,000.
\item Walmart sells 1000 bottles of Coca-Cola for \$1,500. It had previously paid \$1,200 for them.
\item A shipyard builds a cruise ship. It pays wages of \$200,000, interest on loans (from US residents) of \$100,000 and \$300,000
for imported raw materials. The ship is sold for \$1,000,000 to a cruise company. In the same year, the cruise company has revenue
for \$50,000 from operating cruises, pays wages of \$20,000 to its workers and has no other expenses. Half the cruise revenue comes
from tourists who reside in the United States and half comes from tourists who reside abroad.
\item The government collects \$1000 in income taxes from Roger.
\item Roger earns \$4000 for working as a babysitter and pays \$1000 in income taxes.
\end{enumerate}
\resolespaco
\begin{sol}
\begin{notacao}
\item $j$: a producer (a firm, a government unit or a self-employed person). $Q_j\ge0$: its gross output, sales plus the goods it
adds to its inventories (minus the goods it draws from them). $Z_j\ge0$: the intermediate inputs it buys in the period.
\item $VA_j$: value added of $j$. $W_j\ge0$, $R_j\ge0$: the wages and the interest it pays. $\Pi_j$: its profit (proprietor's
income for the self-employed), gross of depreciation; it can be zero or negative.
\item $C$, $G$, $X$, $M\ge0$: consumption, government purchases, exports, imports. $I$: investment, including the change in
inventories, so it can be negative.
\item (G1)--(G4): the accounting identities stated in the first step; the later questions of this list cite them.
\end{notacao}

\passo{method}{For every event, list (i) each producer's value added; (ii) the incomes generated, gross of depreciation;
(iii) the final uses. The three columns rest on four identities:}
\begin{align}
\text{GDP} &= \textstyle\sum_j VA_j, \tag{G1}\label{q1:G1}\\
VA_j &= Q_j-Z_j, \tag{G2}\label{q1:G2}\\
\text{GDP} &= \textstyle\sum_j\big(W_j+R_j+\Pi_j\big),\qquad \Pi_j\equiv VA_j-W_j-R_j, \tag{G3}\label{q1:G3}\\
\text{GDP} &= C+I+G+X-M. \tag{G4}\label{q1:G4}
\end{align}
Transfers and taxes enter none of the three.

\passo{(a) components imported and stored}{Inventories rise by 1, so the firm's output is $Q=1$; the imported input is $Z=1$.
Using \eqref{q1:G2}:}
\begin{align}
VA &= (1)-(1) \nonumber\\
&= 0. \label{q1:aVA}
\end{align}
With no wages or interest, \eqref{q1:G3} makes profit equal to \eqref{q1:aVA}:
\begin{align*}
\Pi &= (0)-(0)-(0)\\
&= 0 .
\end{align*}
The final uses are inventory investment $I=1$ and imports $M=1$; by \eqref{q1:G4}:
\begin{align*}
\text{GDP} &= (1)-(1)\\
&= 0 .
\end{align*}
\begin{center}\footnotesize\begin{tabular}{@{}lr@{\qquad}lr@{\qquad}lr@{}}\toprule
\multicolumn{2}{@{}c}{Production} & \multicolumn{2}{c}{Income} & \multicolumn{2}{c@{}}{Expenditure}\\\midrule
Car manufacturing $1-1$ & 0 & Corporate profits & 0 & Investment (inventories) & 1\\
 & & & & Imports & $-1$\\\midrule
Total & 0 & Total & 0 & Total & 0\\\bottomrule\end{tabular}\end{center}
Nothing was produced in the country: the inventory investment is exactly offset by the import.

\passo{(b) half used, half stored}{Output is the car sold plus the components added to inventory:}
\begin{align}
Q &= (2)+(0.5) \nonumber\\
&= 2.5. \label{q1:bQ}
\end{align}
The intermediate input is the imported 1. Using \eqref{q1:G2} with \eqref{q1:bQ}:
\begin{align}
VA &= (2.5)-(1) \nonumber\\
&= 1.5. \label{q1:bVA}
\end{align}
With no wages or interest, \eqref{q1:G3} gives $\Pi=VA$. Directly, profit is revenue minus the cost of the components used up
(the other 0.5 is still an asset, in inventory):
\begin{align*}
\Pi &= (2)-(0.5)\\
&= 1.5 ,
\end{align*}
which matches \eqref{q1:bVA}. The final uses are $C=2$, inventory investment $I=0.5$ and $M=1$; by \eqref{q1:G4}:
\begin{align*}
\text{GDP} &= (2)+(0.5)-(1)\\
&= 1.5 .
\end{align*}
\begin{center}\footnotesize\begin{tabular}{@{}lr@{\qquad}lr@{\qquad}lr@{}}\toprule
\multicolumn{2}{@{}c}{Production} & \multicolumn{2}{c}{Income} & \multicolumn{2}{c@{}}{Expenditure}\\\midrule
Car manufacturing $2+0.5-1$ & 1.5 & Corporate profits & 1.5 & Consumption & 2\\
 & & & & Investment (inventories) & 0.5\\
 & & & & Imports & $-1$\\\midrule
Total & 1.5 & Total & 1.5 & Total & 1.5\\\bottomrule\end{tabular}\end{center}
Confirmed against Kurlat's own worked Example~1.6 (p.~19), the same event at ten times the scale: 15 by all three approaches.

\passo{(c) army battalion}{Government output has no price, so it is valued at cost, wages plus ammunition:}
\begin{align}
Q_{\mathrm{def}} &= (10{,}000)+(5{,}000) \nonumber\\
&= 15{,}000. \label{q1:cQ}
\end{align}
The ammunition is an intermediate input of defence. Using \eqref{q1:G2} with \eqref{q1:cQ}:
\begin{align}
VA_{\mathrm{def}} &= (15{,}000)-(5{,}000) \nonumber\\
&= 10{,}000. \label{q1:cVAd}
\end{align}
The ammunition maker sells 5{,}000 and uses 2{,}000 of imported materials; by \eqref{q1:G2}:
\begin{align}
VA_{\mathrm{am}} &= (5{,}000)-(2{,}000) \nonumber\\
&= 3{,}000. \label{q1:cVAa}
\end{align}
It pays wages of 1{,}000, so by \eqref{q1:G3} its profit is
\begin{align}
\Pi_{\mathrm{am}} &= (5{,}000)-(2{,}000)-(1{,}000) \nonumber\\
&= 2{,}000. \label{q1:cPi}
\end{align}
Product: summing \eqref{q1:cVAd} and \eqref{q1:cVAa} with \eqref{q1:G1},
\begin{align*}
\text{GDP} &= (10{,}000)+(3{,}000)\\
&= 13{,}000 .
\end{align*}
Income: by \eqref{q1:G3}, the soldiers' wages, the ammunition wages and the profit \eqref{q1:cPi},
\begin{align*}
\text{GDP} &= (10{,}000)+(1{,}000)+(2{,}000)\\
&= 13{,}000 .
\end{align*}
Expenditure: $G$ is the cost \eqref{q1:cQ} and $M=2{,}000$; by \eqref{q1:G4},
\begin{align*}
\text{GDP} &= (15{,}000)-(2{,}000)\\
&= 13{,}000 .
\end{align*}
\begin{center}\footnotesize\begin{tabular}{@{}lr@{\qquad}lr@{\qquad}lr@{}}\toprule
\multicolumn{2}{@{}c}{Production} & \multicolumn{2}{c}{Income} & \multicolumn{2}{c@{}}{Expenditure}\\\midrule
Public defence & 10{,}000 & Wages (soldiers) & 10{,}000 & Government & 15{,}000\\
Manufacturing & 3{,}000 & Wages (ammunition) & 1{,}000 & Imports & $-2{,}000$\\
 & & Corporate profits & 2{,}000 & & \\\midrule
Total & 13{,}000 & Total & 13{,}000 & Total & 13{,}000\\\bottomrule\end{tabular}\end{center}

\passo{(d) Walmart}{Reading: the bottles were bought in an earlier period and sat in Walmart's inventory. Walmart sells 1{,}500
and draws 1{,}200 of bottles from inventory, so its output this year is the retail service:}
\begin{align}
Q_W &= (1{,}500)-(1{,}200) \nonumber\\
&= 300. \label{q1:dQ}
\end{align}
It buys no inputs this year, so by \eqref{q1:G2} with \eqref{q1:dQ}, and then by \eqref{q1:G3} with no wages or interest:
\begin{align*}
VA_W &= (300)-(0)\\
&= 300,\\
\Pi_W &= (300)-(0)-(0)\\
&= 300 .
\end{align*}
The final uses are $C=1{,}500$
and inventory investment $I=-1{,}200$ (the bottles leave inventory); by \eqref{q1:G4}:
\begin{align*}
\text{GDP} &= (1{,}500)+(-1{,}200)\\
&= 300 ,
\end{align*}
equal to \eqref{q1:dQ}.
\begin{center}\footnotesize\begin{tabular}{@{}lr@{\qquad}lr@{\qquad}lr@{}}\toprule
\multicolumn{2}{@{}c}{Production} & \multicolumn{2}{c}{Income} & \multicolumn{2}{c@{}}{Expenditure}\\\midrule
Retail trade & 300 & Corporate profits & 300 & Consumption & 1{,}500\\
 & & & & Investment (inventories) & $-1{,}200$\\\midrule
Total & 300 & Total & 300 & Total & 300\\\bottomrule\end{tabular}\end{center}
If instead the bottles were made by a domestic bottler this same year, the bottler's value added (1{,}200) is also this year's, inventories
do not change and, by \eqref{q1:G1} and \eqref{q1:G4}, GDP becomes
\begin{align*}
\text{GDP} &= (1{,}200)+(300)\\
&= 1{,}500=C .
\end{align*}
State the reading you use.

\passo{(e) shipyard and cruise company}{The shipyard sells the ship for 1{,}000{,}000 and uses 300{,}000 of imported materials;
by \eqref{q1:G2}:}
\begin{align}
VA_S &= (1{,}000{,}000)-(300{,}000) \nonumber\\
&= 700{,}000. \label{q1:eVAs}
\end{align}
It pays wages of 200{,}000 and interest of 100{,}000, so by \eqref{q1:G3}:
\begin{align}
\Pi_S &= (1{,}000{,}000)-(200{,}000)-(100{,}000)-(300{,}000) \nonumber\\
&= 400{,}000. \label{q1:ePis}
\end{align}
The cruise company uses no inputs and pays wages of 20{,}000; by \eqref{q1:G2} and \eqref{q1:G3}:
\begin{align}
VA_{Cr} &= (50{,}000)-(0) \nonumber\\
&= 50{,}000, \label{q1:eVAc}\\
\Pi_{Cr} &= (50{,}000)-(20{,}000) \nonumber\\
&= 30{,}000. \label{q1:ePic}
\end{align}
Product: summing \eqref{q1:eVAs} and \eqref{q1:eVAc} with \eqref{q1:G1},
\begin{align*}
\text{GDP} &= (700{,}000)+(50{,}000)\\
&= 750{,}000 .
\end{align*}
Income: by \eqref{q1:G3}, with the profits \eqref{q1:ePis} and \eqref{q1:ePic},
\begin{align*}
\text{GDP} &= \underbrace{\big((200{,}000)+(20{,}000)\big)}_{\text{wages}}+\underbrace{(100{,}000)}_{\text{interest}}
+\underbrace{\big((400{,}000)+(30{,}000)\big)}_{\text{profits}}\\
&= 220{,}000+100{,}000+430{,}000\\
&= 750{,}000 .
\end{align*}
Expenditure: the ship is an investment good, $I=1{,}000{,}000$; the cruises sold to residents are consumption, $C=25{,}000$, and
those sold to non-residents are exports, $X=25{,}000$; $M=300{,}000$. By \eqref{q1:G4}:
\begin{align*}
\text{GDP} &= (1{,}000{,}000)+(25{,}000)+(25{,}000)-(300{,}000)\\
&= 750{,}000 .
\end{align*}
\begin{center}\footnotesize\begin{tabular}{@{}lr@{\qquad}lr@{\qquad}lr@{}}\toprule
\multicolumn{2}{@{}c}{Production} & \multicolumn{2}{c}{Income} & \multicolumn{2}{c@{}}{Expenditure}\\\midrule
Shipbuilding & 700{,}000 & Wages $200{,}000+20{,}000$ & 220{,}000 & Investment & 1{,}000{,}000\\
Cruise services & 50{,}000 & Interest & 100{,}000 & Consumption & 25{,}000\\
 & & Profits $400{,}000+30{,}000$ & 430{,}000 & Exports & 25{,}000\\
 & & & & Imports & $-300{,}000$\\\midrule
Total & 750{,}000 & Total & 750{,}000 & Total & 750{,}000\\\bottomrule\end{tabular}\end{center}

\passo{(f) income tax}{A transfer from Roger to the government: no good or service is produced. Production 0, Income 0, Expenditure 0.}

\passo{(g) babysitting}{Production: household services with no inputs, so by \eqref{q1:G2}; income: proprietor's income
by \eqref{q1:G3} (Roger is self-employed), with no wages or interest:}
\begin{align*}
VA &= (4{,}000)-(0)\\
&= 4{,}000,\\
\Pi &= (4{,}000)-(0)-(0)\\
&= 4{,}000 .
\end{align*}
The 1{,}000 of tax only splits that income between Roger and the government. Expenditure:
consumption 4{,}000 (the parents buy the service), so \eqref{q1:G4} gives the same total. GDP $=4{,}000$.

\intuicao{The three columns are three views of the same production: who made it, who earned it, who used it. When the columns disagree,
something was double-counted (an intermediate input) or left out (inventories, imports, government output at cost).}
\refbib{Kurlat (2020), \S1.1, pp.~15--22: value added (Ex.~1.2, p.~17), inventories (Ex.~1.3, p.~17), imports (Ex.~1.6, p.~19),
government at cost and transfers (Exs.~1.7--1.8, p.~20). No book answer exists for Exercise 1.1; item (b) is checked against Ex.~1.6.}
\rubrica{
\item 3.5 points per item (a)--(e), 1.25 per item (f)--(g); within an item, the three columns weigh equally.
\item Only one approach given, even with the right total: max.\ a third of the item (the Prova~01 Q1(b) error).
\item Imported input or stored components left out of value added: max.\ half of (a)--(b).
\item Taxes counted as income or expenditure in (f)--(g): 0 for those items.
\item Arithmetic slip with the right layout: $-0.5$ point, once per item.
}
\end{sol}

% -----------------------------------------------------------------------------
\questao{(10 points) \textit{Source: Jones, C.\,I. (2020), \emph{Macroeconomics}, 5th ed., ch.~2, Exercise~2 ``What counts as
GDP (I)?'', p.~39, and Exercise~3 ``What counts as GDP (II)?'', pp.~39--40 --- statements reproduced verbatim, all items.}}
\textit{What this trains:} classifying an event in one line --- new production or not, final or intermediate, which
component of $C+I+G+X-M$ moves --- before writing any table.

\noindent\textbf{2. What counts as GDP (I)?} By how much does GDP rise in each of the following scenarios? Explain.
\begin{enumerate}[label=(\alph*),leftmargin=2.2em,topsep=2pt,itemsep=1pt]
\item You spend \$5,000 on college tuition this semester.
\item You buy a used car from a friend for \$2,500.
\item The government spends \$100 million to build a dam.
\item Foreign graduate students work as teaching assistants at the local university and earn \$5,000 each.
\end{enumerate}
\noindent\textbf{3. What counts as GDP (II)?} By how much does GDP rise in each of the following scenarios? Explain.
\begin{enumerate}[label=(\alph*),leftmargin=2.2em,topsep=2pt,itemsep=1pt]
\item A computer company buys parts from a local distributor for \$1 million, assembles the parts, and sells the resulting
computers for \$2 million.
\item A real estate agent sells a house for \$200,000 that the previous owners had bought 10 years earlier for \$100,000. The agent
earns a commission of \$6,000.
\item During a recession, the government raises unemployment benefits by \$100 million.
\item A new U.S. airline purchases and imports \$50 million worth of airplanes from the European company Airbus.
\item A new European airline purchases \$50 million worth of airplanes from the American company Boeing.
\item A store buys \$100,000 of chocolate from Belgium and sells it to consumers in the United States for \$125,000.
\end{enumerate}
\resolespaco
\begin{sol}
\begin{notacao}
\item $\Delta\text{GDP}$: the change in this year's GDP caused by the event; $\Delta C$, $\Delta I$, $\Delta G$, $\Delta X$, $\Delta M$:
the changes in consumption, investment, government purchases, exports and imports (each of any sign).
\item $VA$: value added, as in \eqref{q1:G1}--\eqref{q1:G2} of Question~1.
\end{notacao}
Taking changes in \eqref{q1:G4}, every item below is read off
\begin{equation}
\Delta\text{GDP}=\Delta C+\Delta I+\Delta G+\Delta X-\Delta M. \label{q2:dY}
\end{equation}

\passo{Exercise 2}{(a) $+5{,}000$: education is a service produced this semester, bought by a household ($C$).
(b) $0$: a used car was produced in an earlier year; selling it transfers an existing asset (a dealer's margin, if any, would count).
(c) $+100$ million: new production of a public structure ($G$; government investment).
(d) $+5{,}000$ per student: the work is done inside the country, so it is domestic product whatever the workers' nationality.}

\passo{Exercise 3}{(a) $+2$ million: the final sale. By value added, \eqref{q1:G1} with the parts treated as intermediate:}
\begin{align*}
\Delta\text{GDP} &= \underbrace{(1\text{ million})}_{\text{computer maker}}+\underbrace{(1\text{ million})}_{\text{parts chain}}\\
&= 2\text{ million}.
\end{align*}
(b) $+6{,}000$: only the agent's service is new production; the house and its capital gain are not.
(c) $0$: unemployment benefits are transfers, not purchases of goods or services.
(d) $0$: $I$ rises by 50 million and $M$ by 50 million; by \eqref{q2:dY},
\begin{align*}
\Delta\text{GDP} &= \Delta I-\Delta M\\
&= (50\text{ million})-(50\text{ million})\\
&= 0 .
\end{align*}
(e) $+50$ million: exports, $\Delta X$ in \eqref{q2:dY}.
(f) $+25{,}000$: $C$ rises by 125{,}000 and $M$ by 100{,}000; by \eqref{q2:dY},
\begin{align*}
\Delta\text{GDP} &= \Delta C-\Delta M\\
&= (125{,}000)-(100{,}000)\\
&= 25{,}000 ,
\end{align*}
which is the store's retail margin, the domestic value added.

\intuicao{GDP counts current production of final goods and services inside the country. Each ``0'' above fails one of those words:
not current (used car, house), not production (transfer, capital gain), not domestic (imported aircraft).}
\refbib{Answers confirmed: Exercise~2 against the worked exercise in Jones (2020), p.~41; Exercise~3 against the Jones
\emph{Instructor's Manual} (2016), ch.~2, Exercise~2 (numbered 3 in the 2020 edition), p.~12 --- same six items, same answers.
Kurlat (2020), Ex.~1.4 (financial investment, p.~18) and Ex.~1.8 (transfers, p.~20).}
\rubrica{
\item 1 point per item (ten items), with the component of $C+I+G+X-M$ named; a right number with no reason: half.
\item Used goods, capital gains or transfers counted in GDP: 0 for that item.
\item Imported aircraft counted as a rise in GDP in 3(d): 0 for that item (it is $I$ and $-M$).
}
\end{sol}

\needspace{12\baselineskip}   % keep each block heading with its first question
% =============================================================================
\bloco{2}{Harder book items and conceptual classification}{Kurlat \S\S1.1, 2.2, 4.3, 5.2, 5.4; Jones ch.~2 \,\textbullet\, 35 points}

\questao{(10 points) \textit{Source: Kurlat, P. (2020), \emph{A Course in Modern Macroeconomics}, Exercise 5.8 ``National Accounts and
The Golden Rule'', pp.~97--98 --- statement reproduced verbatim, all items (a)--(d).}}
\par\smallskip\noindent\fbox{\parbox{\dimexpr\linewidth-2\fboxsep-2\fboxrule}{\small\textbf{Context from the book (added by this list, from
Kurlat \S1.1, p.~21; \S\S4.1--4.4, pp.~55--67; \S\S5.1--5.2, pp.~76--77):} \emph{Accounts:} the conventions in the box of Question~1, plus
depreciation: machines and buildings lose value over time; GDP is \emph{gross} because it is measured before depreciation, so when a firm
subtracts depreciation to compute its profits, the income approach adds it back as a separate line (Ex.~1.10, p.~21). The savings rate is
$S/Y=(Y-C)/Y$, savings as a fraction of income, and in a closed economy savings equal investment, $S=I$ (p.~56).
\emph{Solow model:} with $Y$ output, $K$ capital, $L$ labor (everyone works, so $L$ is also the population), the Cobb--Douglas
production function is $Y=K^{\alpha}L^{1-\alpha}$, $\alpha\in(0,1)$ \,(4.1.2), p.~55 (writing it $AK^{\alpha}L^{1-\alpha}$, with $A$ the
level of technology, ``doesn't make much difference'', p.~77). The population grows at rate $n$, $L_{t+1}=(1+n)L_t$ (p.~55); people
save a constant fraction $s$ of income, so $I=sY$ \,(4.1.3); every piece of capital loses a fraction $\delta$ of its value each period,
so $K_{t+1}=(1-\delta)K_t+I_t$ \,(4.1.4) (pp.~56--57). Per worker, $y\equiv Y/L$, $k\equiv K/L$ and $y=f(k)\equiv F(K/L,1)$ \,(4.2.1)
(p.~57). In a steady state ($ss$), $sf(k_{ss})=(\delta+n)k_{ss}$ \,(4.3.2) and consumption is $c_{ss}=(1-s)y_{ss}$ \,(4.3.1); the
\emph{Golden Rule} is the saving behavior that maximizes steady-state consumption, and the Golden Rule capital stock $k_{gr}$ satisfies
$f'(k_{gr})=\delta+n$ \,(4.3.4) (pp.~62--63). \emph{Factor markets} are perfectly competitive, so the rental rate of capital and the wage
equal the marginal products, $r^K=F_K(K,L)$ \,(4.4.8) and $w=F_L(K,L)$ \,(4.4.9) (p.~66), and all output is paid to workers or to
capital owners, $r^KK+wL=F(K,L)$ (Prop.~4.2, p.~67). With Cobb--Douglas the labor share is $wL/Y=1-\alpha$ \,(5.1.1), so the capital
share is $\alpha$ (p.~76).}}
\par\smallskip
\textit{What this trains:} the three approaches with an investment good and depreciation, then reading factor shares and
the saving rate off the accounts and using them in the golden rule --- Prova~01 Q1 and Q2(b) in one exercise.

\noindent In 2018, the following events happened in Aurum, which is a closed economy that uses the Denarius (plural: Denarii) as its currency.
\begin{itemize}[leftmargin=2.2em,topsep=2pt,itemsep=1pt]
\item Grano, Inc.\ hired Cornelia to plant and harvest 500 tons of wheat on a large plot of land using a tractor. Both the land and
the tractor are owned by Grano, Inc.
  \begin{itemize}[label=--,topsep=0pt,itemsep=0pt]
  \item It paid Cornelia 200 Denarii for her work on the harvest.
  \item It sold the wheat to Panem, Inc.\ for 1,000 Denarii.
  \end{itemize}
\item The tractor the Cornelia used:
  \begin{itemize}[label=--,topsep=0pt,itemsep=0pt]
  \item Was new at the beginning of the year. It was worth 2,000 Denarii.
  \item Was no longer new at the end of the year. By then it was worth 1,800 Denarii.
  \end{itemize}
\item Panem, Inc.\ hired Gaius to grind the wheat in its mill and produce a million loaves of bread.
  \begin{itemize}[label=--,topsep=0pt,itemsep=0pt]
  \item It paid Gaius 1,400 Denarii for his work.
  \item It sold the bread for 4,000 Denarii to the hungry citizens, who eat it.
  \end{itemize}
\item Fabrica, Inc.\ hired Livia to build a combine harvester.
  \begin{itemize}[label=--,topsep=0pt,itemsep=0pt]
  \item It paid Livia 400 Denarii for her work.
  \item It sold the combine harvester to Grano, Inc.\ for 1,000 Denarii.
  \end{itemize}
\end{itemize}
\begin{enumerate}[label=(\alph*),leftmargin=2.2em,topsep=2pt,itemsep=1pt]
\item Construct GDP accounts (in Denarii) for Aurum by production, income and expenditure.
\item What were the income shares of labor and capital?
\item What was the savings rate?
\item Assume the following:
  \begin{itemize}[topsep=0pt,itemsep=0pt]
  \item the savings rate is constant,
  \item the production function has the form:
  \[
  Y=AK^\alpha L^{1-\alpha}
  \]
  and there is no technological progress,
  \item the depreciation rate for the tractor is typical for the capital stock as a whole,
  \item the population is constant.
  \end{itemize}
In the long run, would the level of consumption be higher if the economy slightly increased its saving rate?
\end{enumerate}
\resolespaco
\begin{sol}
\begin{notacao}
\item Grano (G), Panem (P), Fabrica (F): the three firms. $VA_j$, $\Pi_j$ and the identities \eqref{q1:G1}--\eqref{q1:G4}: as in Question~1.
\item $D\ge0$: depreciation of Grano's tractor. $S$: national saving. $Y>0$: GDP; $K>0$: the capital stock; $L>0$: labour.
\item $A>0$: total factor productivity. $\alpha\in(0,1)$: the capital exponent of $Y=AK^{\alpha}L^{1-\alpha}$, equal to the capital
share under competitive factor pay. $\delta\in(0,1)$: depreciation rate. $n\ge0$: labour-force growth (here $n=0$).
\item $s\in(0,1)$: saving rate. $k\equiv K/L$, $y\equiv Y/L$, $c$: capital, output and consumption per worker; the subscript $ss$
marks the steady state. $s_{GR}$: the golden-rule saving rate. $MPK$: marginal product of capital.
\end{notacao}

\passo{(a) accounts}{Grano uses no intermediate inputs (the tractor and the land are capital); the wheat is intermediate for Panem;
the combine harvester is a final good, investment by Grano. By \eqref{q1:G2}:}
\begin{align}
VA_G &= (1{,}000)-(0) \nonumber\\
&= 1{,}000, \label{q3:vaG}\\
VA_P &= (4{,}000)-(1{,}000) \nonumber\\
&= 3{,}000, \label{q3:vaP}\\
VA_F &= (1{,}000)-(0) \nonumber\\
&= 1{,}000. \label{q3:vaF}
\end{align}
Summing \eqref{q3:vaG}--\eqref{q3:vaF} with \eqref{q1:G1}:
\begin{align}
Y &= (1{,}000)+(3{,}000)+(1{,}000) \nonumber\\
&= 5{,}000. \label{q3:Y}
\end{align}
The tractor loses value from 2{,}000 to 1{,}800:
\begin{align}
D &= (2{,}000)-(1{,}800) \nonumber\\
&= 200. \label{q3:D}
\end{align}
Capital income by \eqref{q1:G3}, value added minus wages:
\begin{align*}
\Pi_G &= (1{,}000)-(200)\\
&= 800,\\
\Pi_P &= (4{,}000)-(1{,}000)-(1{,}400)\\
&= 1{,}600,\\
\Pi_F &= (1{,}000)-(400)\\
&= 600.
\end{align*}
Grano's 800 splits into the depreciation \eqref{q3:D} and a profit of
\begin{align*}
\Pi_G-D &= (800)-(200)\\
&= 600 .
\end{align*}
Adding up wages and profits:
\begin{align}
\text{wages} &= (200)+(1{,}400)+(400) \nonumber\\
&= 2{,}000, \label{q3:W}\\
\text{profits} &= (600)+(1{,}600)+(600) \nonumber\\
&= 2{,}800. \label{q3:Pi}
\end{align}
Income: by \eqref{q1:G3}, adding \eqref{q3:W}, \eqref{q3:Pi} and \eqref{q3:D},
\begin{align*}
Y &= (2{,}000)+(2{,}800)+(200)\\
&= 5{,}000 .
\end{align*}
Expenditure: by \eqref{q1:G4}, with $C=4{,}000$ (Panem's food) and $I=1{,}000$ (the harvester),
\begin{align*}
Y &= (4{,}000)+(1{,}000)\\
&= 5{,}000 .
\end{align*}
All three equal \eqref{q3:Y}.
\begin{center}\footnotesize\begin{tabular}{@{}lr@{\qquad}lr@{\qquad}lr@{}}\toprule
\multicolumn{2}{@{}c}{Production} & \multicolumn{2}{c}{Income} & \multicolumn{2}{c@{}}{Expenditure}\\\midrule
Agriculture (Grano) & 1{,}000 & Wages $200+1{,}400+400$ & 2{,}000 & Consumption & 4{,}000\\
Food (Panem) & 3{,}000 & Profits $600+1{,}600+600$ & 2{,}800 & Investment & 1{,}000\\
Machinery (Fabrica) & 1{,}000 & Depreciation & 200 & & \\\midrule
Total & 5{,}000 & Total & 5{,}000 & Total & 5{,}000\\\bottomrule\end{tabular}\end{center}
\passo{(b) shares}{Dividing \eqref{q3:W} by \eqref{q3:Y}:}
\begin{align*}
\text{labor share} &= \frac{(2{,}000)}{(5{,}000)}\\
&= 0.4 .
\end{align*}
Dividing the sum of \eqref{q3:Pi} and \eqref{q3:D} by \eqref{q3:Y} (gross of depreciation, as GDP is gross):
\begin{align}
\text{capital share} &= \frac{(2{,}800)+(200)}{(5{,}000)} \nonumber\\
&= 0.6. \label{q3:ksh}
\end{align}
\passo{(c) saving rate}{Closed economy, no government, so saving is output not consumed, and it equals investment:}
\begin{align*}
S &= Y-C\\
&= (5{,}000)-(4{,}000)\\
&= 1{,}000=I .
\end{align*}
Dividing by \eqref{q3:Y}:
\begin{align}
s &= \frac{(1{,}000)}{(5{,}000)} \nonumber\\
&= 0.2. \label{q3:s}
\end{align}
\passo{(d) golden rule}{With $Y=AK^\alpha L^{1-\alpha}$, competitive factor pay makes the capital share equal to $\alpha$, so by
\eqref{q3:ksh} $\alpha=0.6$. Dividing \eqref{q3:D} by the tractor's value:}
\begin{align}
\delta &= \frac{(200)}{(2{,}000)} \nonumber\\
&= 0.10, \label{q3:delta}
\end{align}
and $n=0$. In steady state saving replaces depreciation, $sY=\delta K$; per worker, substituting $y=Ak^{\alpha}$:
\begin{align}
sy &= \delta k, \nonumber\\
s\big(Ak^{\alpha}\big) &= \delta k. \label{q3:ss}
\end{align}
Divide \eqref{q3:ss} by $\delta k^{\alpha}$ and raise to the power $1/(1-\alpha)$:
\begin{align}
k^{1-\alpha} &= \frac{sA}{\delta}, \nonumber\\
k_{ss} &= \Big(\frac{sA}{\delta}\Big)^{\frac{1}{1-\alpha}}. \label{q3:kss}
\end{align}
Substitute \eqref{q3:kss} into $y=Ak^{\alpha}$:
\begin{align}
y_{ss} &= A\bigg(\Big(\frac{sA}{\delta}\Big)^{\frac{1}{1-\alpha}}\bigg)^{\alpha} \nonumber\\
&= A\Big(\frac{sA}{\delta}\Big)^{\frac{\alpha}{1-\alpha}} \nonumber\\
&= A^{\frac{1}{1-\alpha}}\Big(\frac{s}{\delta}\Big)^{\frac{\alpha}{1-\alpha}}. \label{q3:yss}
\end{align}
Consumption is the part of output not saved:
\begin{equation}
c_{ss}=(1-s)\,y_{ss}. \label{q3:css}
\end{equation}
Substitute \eqref{q3:yss} into \eqref{q3:css}, take logs and differentiate with respect to $s$:
\begin{align}
c_{ss} &= (1-s)\Big(A^{\frac{1}{1-\alpha}}\Big(\frac{s}{\delta}\Big)^{\frac{\alpha}{1-\alpha}}\Big), \nonumber\\
\ln c_{ss} &= \ln(1-s)+\frac{1}{1-\alpha}\ln A+\frac{\alpha}{1-\alpha}\big(\ln s-\ln\delta\big), \nonumber\\
\frac{d\ln c_{ss}}{ds} &= -\frac{1}{1-s}+\frac{\alpha}{(1-\alpha)s}. \label{q3:dlnc}
\end{align}
Evaluate \eqref{q3:dlnc} at $\alpha=0.6$ and \eqref{q3:s}:
\begin{align*}
\frac{d\ln c_{ss}}{ds} &= -\frac{1}{1-(0.2)}+\frac{(0.6)}{\big(1-(0.6)\big)(0.2)}\\
&= -\frac{1}{0.8}+\frac{0.6}{0.4\times0.2}\\
&= -1.25+7.5\\
&= 6.25>0 .
\end{align*}
Multiplying \eqref{q3:dlnc} by $(1-\alpha)s(1-s)>0$ leaves $\alpha(1-s)-(1-\alpha)s=\alpha-s$, so the sign is positive exactly
when $s<\alpha$:
\begin{equation}
s_{GR}=\alpha=0.6. \label{q3:sgr}
\end{equation}
Cross-check with Kurlat's condition $f'(k_{gr})=\delta+n$. From \eqref{q3:ss}, with \eqref{q3:s} and \eqref{q3:delta}:
\begin{align*}
\frac{K}{Y} &= \frac{s}{\delta}\\
&= \frac{(0.2)}{(0.1)}\\
&= 2 .
\end{align*}
So, by the marginal product of the Cobb--Douglas function,
\begin{align*}
MPK &= \alpha\frac{Y}{K}\\
&= \frac{(0.6)}{(2)}\\
&= 0.3>\delta+n=0.1 ;
\end{align*}
capital is below the golden rule. \textbf{Yes:} a slightly higher saving rate raises long-run consumption, by \eqref{q3:sgr}.
\intuicao{The national accounts hand you the two numbers the golden rule needs: the capital share ($\alpha$) and the saving rate ($I/Y$).
Saving 20\% when capital earns 60\% of income means the marginal unit of capital pays far more than it costs to maintain.}
\refbib{Kurlat (2020), \S1.1, Ex.~1.10 (depreciation added back), p.~21; \S4.3, eq.~(4.3.4), p.~63; \S5.2, p.~78. No book answer exists;
every figure is verified by the check script.}
\rubrica{
\item (a) 4 points: each column 1.3; the combine harvester counted as intermediate (GDP 4{,}000): max.\ 1.5 for (a).
\item (b) 2 points; (c) 1 point.
\item (d) 3 points: $s_{GR}=\alpha$ derived or justified 1.5, numbers and conclusion 1.5. Comparing $s$ with the labor share: max.\ 1.
}
\end{sol}

% -----------------------------------------------------------------------------
\needspace{22\baselineskip}   % keep the source line with its context box
\questao{(12 points) \textit{Source: Kurlat, P. (2020), \emph{A Course in Modern Macroeconomics}, Exercise 5.7 ``GDP Accounting,
Interest Rates and Growth Accounting'', pp.~96--97 --- statement reproduced verbatim, all items (a)--(e).}}
\par\smallskip\noindent\fbox{\parbox{\dimexpr\linewidth-2\fboxsep-2\fboxrule}{\small\textbf{Context from the book (added by this list, from
Kurlat \S4.4, pp.~66 and 69; \S5.4, pp.~86--87):} \emph{Accounts:} imports, government output at cost and transfers as in the box of
Question~1; depreciation, the savings rate, the Cobb--Douglas function and its factor shares (labor $1-\alpha$, capital $\alpha$) as in
the box of Question~3. $F_K$ is the marginal product of capital, $\partial Y/\partial K$. \emph{Interest rate:} in competitive markets
the rental rate of capital is $r^K=F_K(K,L)$ \,(4.4.8) (p.~66); someone who borrows at the real interest rate $r$ to invest in capital
earns $r^K-\delta$, so the loan market clears only if
\[ r_{t+1}=r^K_{t+1}-\delta \tag{4.4.12} \]
(p.~69), with $\delta$ the rate at which capital loses value. \emph{Growth accounting:} with $Y_t=F(K_t,L_t,A_t)$ \,(5.4.1), $A$
technology, and GDP growth $g_Y\equiv Y_{t+1}/Y_t-1$, a first-order Taylor approximation gives (5.4.2): $g_Y\approx$ capital share
$\times$ growth rate of capital $+$ labor share $\times$ growth rate of labor $+$ Solow residual, the shares being $F_KK/F$ and $F_LL/F$
(pp.~86--87). The residual is measured as ``Solow Residual $=$ GDP growth $-$ Capital share $\times$ Capital Growth $-$ Labor Share
$\times$ Labor Growth'' and is also called Total Factor Productivity, TFP (p.~87).}}
\par\smallskip
\textit{What this trains:} an open economy with imported capital, imported consumption, exports, government output at cost and
income taxes; then the same accounts feed the interest rate and a Solow residual.

\noindent Proletaria is a mythical country in Central Asia, where the currency is the ruble. The exchange rate is 10 rubles per euro.

Suppose the production function in Proletaria is given by:
\[
Y=K^\alpha L^{1-\alpha}
\]
\begin{enumerate}[label=(\alph*),leftmargin=2.2em,topsep=2pt,itemsep=1pt]
\item Find expressions for the output-to-capital ratio $\frac YK$ and the marginal product of capital $F_K$ as functions of $K$
and $L$.
\end{enumerate}
During 2014, the following events took place in Proletaria.
\begin{itemize}[leftmargin=2.2em,topsep=2pt,itemsep=1pt]
\item Kapitas, the main manufacturer in the country, imported an industrial welder made in Germany, for which it paid 1,000 euro.
\item By the end of the year, the industrial welder was no longer new. Its estimated value in the resale market was 960 euro.
\item 100 workers worked for Kapitas the entire year making screws and nails, using the new welder. Each of them received wages
for 350 rubles.
\item Each of the Kapitas workers paid 100 rubles in income taxes.
\item The total output of Kapitas consisted of 1 million screws and 1 million nails. All of it was exported to Austria for a total
of 4,000 euro.
\item The government of Proletaria employed one of the 100 workers as Chief of the Secret Police (in addition to his factory job)
to maintain law and order, and paid her a salary of 10,000 rubles.
\item The workers ate beef imported from France, which cost a total of 2,500 euro.
\end{itemize}
\begin{enumerate}[label=(\alph*),leftmargin=2.2em,topsep=2pt,itemsep=1pt,start=2]
\item Construct GDP accounts (in rubles) for Proletaria by production, income and expenditure.
\item What were the income shares of labor and capital?
\item Suppose that we know that the capital stock in Proletaria is 100,000 rubles and that the depreciation rate of the
industrial welder is typical for this country's capital stock. What interest rate should we expect to observe?
\item During 2015, additional investment in Proletaria has exceeded depreciation so that the capital stock now stands at
120,000 rubles. The labor force also grew thanks to immigration, and now consists of 105 workers instead of 100. GDP during 2015
was 55,000 rubles (prices were constant). How much of the growth in GDP between 2014 and 2015 can be attributed to growth in
Total Factor Productivity?
\end{enumerate}
\resolespaco
\begin{sol}
\begin{notacao}
\item $Y>0$: GDP; $K>0$: the capital stock; $L>0$: labour (workers); all in the exercise's units, values in rubles.
\item $\alpha\in(0,1)$: the capital exponent of $Y=K^{\alpha}L^{1-\alpha}$, equal to the capital share under competition.
$F_K$: marginal product of capital. $\delta\in(0,1)$: depreciation rate. $r$: the real interest rate.
\item $g_Y$, $g_K$, $g_L$: growth rates of $Y$, $K$, $L$; the Solow residual is the growth of TFP.
\item $C$, $I$, $G$, $X$, $M$, $\Pi$ and the identities \eqref{q1:G1}--\eqref{q1:G4}: as in Question~1. The exchange rate is
10 rubles per euro.
\end{notacao}

\passo{(a) ratios}{Write the technology as}
\begin{equation}
Y=K^{\alpha}L^{1-\alpha}. \label{q4:Y}
\end{equation}
Divide \eqref{q4:Y} by $K$:
\begin{align*}
\frac YK &= K^{\alpha-1}L^{1-\alpha}\\
&= \Big(\frac KL\Big)^{\alpha-1}.
\end{align*}
Differentiate \eqref{q4:Y} with respect to $K$ and use \eqref{q4:Y} again:
\begin{align}
F_K &= \alpha K^{\alpha-1}L^{1-\alpha} \nonumber\\
&= \alpha\frac{\big(K^{\alpha}L^{1-\alpha}\big)}{K} \nonumber\\
&= \alpha\frac YK. \label{q4:FK}
\end{align}
\passo{(b) accounts, in rubles (10 rubles per euro)}{Convert each euro amount, and compute the wage bill and the taxes:}
\begin{align*}
\text{welder (imported investment)} &= (1{,}000)\times(10)\\
&= 10{,}000,\\
\text{exports} &= (4{,}000)\times(10)\\
&= 40{,}000,\\
\text{beef (imported consumption)} &= (2{,}500)\times(10)\\
&= 25{,}000,\\
\text{Kapitas wages} &= (100)\times(350)\\
&= 35{,}000,\\
\text{income taxes} &= (100)\times(100)\\
&= 10{,}000.
\end{align*}
Kapitas's gross profit, by \eqref{q1:G3}:
\begin{align*}
\Pi &= (40{,}000)-(35{,}000)\\
&= 5{,}000 .
\end{align*}
It splits into depreciation and net profit:
\begin{align*}
\text{depreciation} &= \big((1{,}000)-(960)\big)\times(10)\\
&= 400,\\
\text{net profit} &= (5{,}000)-(400)\\
&= 4{,}600.
\end{align*}
The police chief's salary, 10{,}000, is government output valued at cost. The income taxes are a transfer and appear nowhere.
Product, by \eqref{q1:G1}--\eqref{q1:G2}:
\begin{align}
Y &= (40{,}000)+(10{,}000) \nonumber\\
&= 50{,}000. \label{q4:Y0}
\end{align}
Income, by \eqref{q1:G3}:
\begin{align*}
Y &= (35{,}000)+(10{,}000)+(5{,}000)\\
&= 50{,}000 .
\end{align*}
Imports are the welder plus the beef; expenditure, by \eqref{q1:G4}:
\begin{align*}
M &= (10{,}000)+(25{,}000)\\
&= 35{,}000,\\
Y &= (25{,}000)+(10{,}000)+(10{,}000)+(40{,}000)-(35{,}000)\\
&= 50{,}000 .
\end{align*}
\begin{center}\footnotesize\begin{tabular}{@{}lr@{\qquad}lr@{\qquad}lr@{}}\toprule
\multicolumn{2}{@{}c}{Production} & \multicolumn{2}{c}{Income} & \multicolumn{2}{c@{}}{Expenditure}\\\midrule
Manufacturing & 40{,}000 & Wages (Kapitas) & 35{,}000 & Consumption & 25{,}000\\
Public administration & 10{,}000 & Wages (government) & 10{,}000 & Investment & 10{,}000\\
 & & Profit 4{,}600 $+$ deprec.\ 400 & 5{,}000 & Government & 10{,}000\\
 & & & & Exports & 40{,}000\\
 & & & & Imports (welder $+$ beef) & $-35{,}000$\\\midrule
Total & 50{,}000 & Total & 50{,}000 & Total & 50{,}000\\\bottomrule\end{tabular}\end{center}
\passo{(c) shares}{Dividing all wages by \eqref{q4:Y0}:}
\begin{align*}
\text{labor share} &= \frac{(35{,}000)+(10{,}000)}{(50{,}000)}\\
&= 0.9 .
\end{align*}
Dividing the gross profit by \eqref{q4:Y0}:
\begin{align}
\alpha=\text{capital share} &= \frac{(5{,}000)}{(50{,}000)} \nonumber\\
&= 0.1. \label{q4:alpha}
\end{align}
\passo{(d) interest rate}{The depreciation rate is the value lost (the same 40 euros behind the 400 rubles of (b)) over the
value at the start:}
\begin{align}
\delta &= \frac{(1{,}000)-(960)}{(1{,}000)} \nonumber\\
&= 0.04. \label{q4:delta}
\end{align}
The interest rate is the marginal product of capital net of depreciation; substitute \eqref{q4:FK}:
\begin{align}
r &= F_K-\delta \nonumber\\
&= \Big(\alpha\frac YK\Big)-\delta. \label{q4:r}
\end{align}
Substitute \eqref{q4:alpha}, \eqref{q4:Y0}, $K=100{,}000$ and \eqref{q4:delta} into \eqref{q4:r}:
\begin{align*}
r &= (0.1)\times\frac{(50{,}000)}{(100{,}000)}-(0.04)\\
&= 0.05-0.04\\
&= 0.01,
\end{align*}
i.e.\ 1\%.
\passo{(e) growth accounting}{Growth rates from one year to the next:}
\begin{align*}
g_Y &= \frac{(55{,}000)}{(50{,}000)}-1\\
&= 10\%,\\
g_K &= \frac{(120{,}000)}{(100{,}000)}-1\\
&= 20\%,\\
g_L &= \frac{(105)}{(100)}-1\\
&= 5\%.
\end{align*}
The Solow residual is output growth not explained by factor growth, weighted by the shares:
\begin{equation}
\text{residual}=g_Y-\alpha\,g_K-(1-\alpha)\,g_L. \label{q4:res}
\end{equation}
Substitute the growth rates and \eqref{q4:alpha} into \eqref{q4:res}:
\begin{align*}
\text{residual} &= (10\%)-(0.1)\times(20\%)-\big(1-(0.1)\big)\times(5\%)\\
&= 10\%-0.1\times20\%-0.9\times5\%\\
&= 10\%-2\%-4.5\%\\
&= 3.5\%.
\end{align*}
So 3.5 of the 10 percentage points, 35\% of the growth, is TFP. With exact logs in \eqref{q4:res}:
\begin{align*}
\text{residual} &= (\ln1.1)-(0.1)(\ln1.2)-(0.9)(\ln1.05)\\
&= 0.0953-0.0182-0.0439\\
&= 0.0332,
\end{align*}
34.8\% of $\ln1.1$.
\intuicao{Imports are subtracted because $C$ and $I$ include goods made abroad; the welder raises investment and imports by the same
amount and leaves GDP untouched. The income side is what gives $\alpha$, and $\alpha$ is what turns GDP into $r$ and into TFP.}
\refbib{Kurlat (2020), \S1.1, Exs.~1.6--1.8, pp.~19--20; \S5.2, $r=F_K-\delta$, p.~78; \S5.4, eq.~(5.4.2), p.~87. No book answer exists.}
\rubrica{
\item (a) 1 point; (b) 5 points (each column 1.5; conversion at 10 rubles per euro 0.5).
\item Income taxes counted in GDP, or the welder counted without its import: max.\ 2 for (b).
\item (c) 2 points (government wages are labor income; using $35{,}000/40{,}000$ with a stated reason: 1.5).
\item (d) 2 points: $r$ without $-\delta$: max.\ 1. (e) 2 points.
}
\end{sol}

% -----------------------------------------------------------------------------
\questao{(10 points) Conceptual --- answer in words, no numbers. For each event, say whether it adds to this year's GDP;
if it does, name the producer whose value added rises, the income it generates and the final use ($C$, $I$, $G$ or $X$) it
enters; if it does not, say why. Two to four sentences each.}
\noindent\textbf{a)} A search engine is free for households and is paid for by advertising sold to firms.\\
\textbf{b)} A bank's own IT staff write a risk-management program the bank will use for years; a drug company's
laboratory runs a research project for its own future products.\\
\textbf{c)} A household sells its five-year-old car to another household through a used-car dealer.\\
\textbf{d)} One family owns the house it lives in; next door an identical house is rented to a tenant.\\
\textbf{e)} An accounting firm prepares one tax return for a manufacturing company and one for a household.\\
\textbf{f)} The government pays retirement pensions and pays the salaries of public-school teachers.
\resolespaco
\begin{sol}
\begin{notacao}
\item $C$, $I$, $G$, $X$: the final uses of \eqref{q1:G4} (consumption, investment, government purchases, exports).
\item Value added, income and the three approaches: \eqref{q1:G1}--\eqref{q1:G3} of Question~1. No numbers are needed.
\end{notacao}
\passo{a)}{The ads are a service bought by firms, so they are an \emph{intermediate} input: the network's value added is counted, but no
final use appears for it; it is embedded in the price of the advertisers' goods. The search itself has a zero price and is no market
transaction, so it adds nothing beyond that. The users' gain is real welfare that GDP does not measure.}
\passo{b)}{Both are \emph{gross fixed capital formation} in intellectual property products: durable assets that yield services for several
years. With no market price, each is valued at its production cost. Production: the bank's and the company's output includes the
asset; income: the wages paid plus an operating surplus; expenditure: $I$. It is not consumption (no household buys it) and not an
intermediate input (it is not used up within the year).}
\passo{c)}{The car was produced, and counted, five years ago; changing its owner is not new production. Only the dealer's service is
current output: its margin is the dealer's value added, its profit and wages, and the buyer's consumption of a trade service.}
\passo{d)}{Both houses produce housing services that enter $C$. For the tenant, the rent is paid in the market; for the owner, an
\emph{imputed} rent is counted as both consumption and income of the owner. Without it, GDP would fall every time a tenant bought the
house she lives in, with no change in what is produced.}
\passo{e)}{The same service is \emph{intermediate} when the company buys it (a cost of making its goods, already inside their price) and
\emph{final} when the household buys it ($C$). The accounting firm's value added counts in both cases; only the final-use column
differs. The classification follows the buyer and the use, never the product.}
\passo{f)}{Pensions are transfers: no good or service is produced in exchange, so they enter neither $G$ nor GDP. Teachers produce
education with no market price; it is valued at cost: value added of public education $=$ the wages $=$ income, and the final use is
$G$.}
\intuicao{One question settles every case: was something produced this year, and who used it up? Firms using it $\Rightarrow$
intermediate; households $\Rightarrow C$; firms keeping it for years $\Rightarrow I$; government $\Rightarrow G$; nothing produced $\Rightarrow$ 0.}
\refbib{Kurlat (2020), \S1.1: inventories and investment (Ex.~1.3, p.~17), financial investment (Ex.~1.4, p.~18), durables and imputed rent
(Ex.~1.5, pp.~18--19), government at cost and transfers (Exs.~1.7--1.8, p.~20), non-market activity (Ex.~1.11, p.~22); \S2.2, p.~33.
Jones (2020), ch.~2, p.~22: intellectual property products in $I$; government purchases vs.\ transfers.}
\rubrica{
\item 10/6 points per item, judged on the classification and the reason.
\item Own-account software or R\&D called consumption (the Prova~01 Q1(c) error), or the free service counted as a final use: 0 for that item.
\item Right classification with no reason: half the item.
}
\end{sol}

% -----------------------------------------------------------------------------
\needspace{28\baselineskip}   % keep the source line with its context box
\questao{(3 points) \textit{Source: Kurlat, P. (2020), \emph{A Course in Modern Macroeconomics}, Exercise 1.4 ``Drugs and
Prostitution'', p.~29, and Exercise 2.4 ``Construction and the Armed Forces'', p.~41; Jones, C.\,I. (2020), \emph{Macroeconomics},
5th ed., ch.~2, Exercise~9 ``Earthquakes and GDP'', p.~41 --- statements reproduced verbatim; the closing instruction is added
by this list.}}
\par\smallskip\noindent\fbox{\parbox{\dimexpr\linewidth-2\fboxsep-2\fboxrule}{\small\textbf{Context from the books (added by this list,
from Kurlat \S1.1, pp.~15 and 21; Kurlat \S2.2, pp.~33--34; Jones ch.~2, pp.~19, 22 and 26):} \emph{Kurlat 1.4.} GDP records the value
of everything produced in the country in a year (p.~15); activity that involves no market transaction is usually not included: ``Most
of the time, however, we compute the value of an activity only if it is sold in the market'' (p.~21). The article is not reproduced here;
its address names its subject: Italy's inclusion of illicit activities in its GDP figures.
\emph{Kurlat 2.4.} Jones and Klenow measure the living standard of a country, Utilia, by the number $\lambda$ that makes a person
indifferent between living a year as a random resident of Utilia and as a random resident of the US with everyone's consumption
multiplied by $\lambda$ (p.~33), using
\[ u(c,l,a)=\mathbb{E}\Big[\Big(\bar u+\frac{c^{1-\sigma}}{1-\sigma}-\theta(1-l)^2\Big)a\Big], \tag{2.2.1} \]
\[ u\big(\lambda c^{US},l^{US},a^{US}\big)=u\big(c^{Utilia},l^{Utilia},a^{Utilia}\big), \tag{2.2.2} \]
where $c$ is consumption, $l$ leisure (the fraction of time spent on activities not counted in GDP), $a$ equals 1 if alive and 0 if dead,
and $\bar u$, $\sigma$, $\theta$ are parameters (pp.~33--34). What enters (2.2.1) is consumption, not production: output dedicated to
investment or exported gives no utility during the year (investment benefits future residents), and public and private consumption are
assumed to give the same utility, so what enters is $C+G$ of (1.1.1) (box of Question~1), more precisely $C$ plus the government-consumption part of $G$
(p.~34). \emph{Jones 9.} GDP is ``the market value of the final goods and services produced in an economy over a certain period''
(p.~19); the construction of new homes is part of investment (p.~22); GDP counts only new production (p.~26).}}
\par\smallskip
\textit{What this trains:} separating what GDP measures (market production) from what it does not (legality, welfare, the
stock of wealth).

\noindent\textbf{Kurlat 1.4 Drugs and Prostitution.} Read the following article:
\texttt{http://\allowbreak www.economist.com/\allowbreak news/\allowbreak finance-\allowbreak and-\allowbreak economics/\allowbreak
21603073-\allowbreak italys-\allowbreak inclusion-\allowbreak illicit-\allowbreak activities-\allowbreak its-\allowbreak figures-\allowbreak
excites-\allowbreak much-\allowbreak interest-\allowbreak sex}. What do you think? Should drug
production and prostitution be included in the calculation of GDP?

\noindent\textbf{Kurlat 2.4 Construction and the Armed Forces.} Suppose we compared welfare in two otherwise identical countries.
Country A dedicates 50\% of GDP to maintaining a large army. Country B dedicates 50\% of GDP to building houses. In which of the
two countries would welfare be higher, according to the measure used by Jones and Klenow. What do you think of this?

\noindent\textbf{Jones 9. Earthquakes and GDP:} Suppose the rural part of a country is hit by a major earthquake that destroys
10 percent of the country's housing stock. The government and private sector respond with a major construction effort to help
rebuild houses. Discuss how this episode is likely to affect (a) the economic well-being of the people in the country and (b) the
economy's measured GDP.

\noindent\textit{Added (not in the book):} answer each of the three in words.
\resolespaco
\begin{sol}
\begin{notacao}
\item $C$, $G$, $I$: consumption, government purchases and investment, as in \eqref{q1:G4}.
\item Jones--Klenow measure: Kurlat's welfare index, built on consumption $C+G$ (Kurlat 2020, \S2.2).
\end{notacao}
\passo{Kurlat 1.4}{GDP measures market production, not its moral worth; drugs and prostitution are voluntary transactions that pay
incomes and meet final demand, so the accounts' logic says include them. Including them also keeps GDP comparable across
countries with different laws and over time when laws change (legalisation would otherwise raise measured GDP with no change in
activity). Against: the figures are estimates built on strong assumptions. Theft, by contrast, is excluded: it is a forced transfer,
not production (Williamson 2014, solutions manual, ch.~2, problem~9).}
\passo{Kurlat 2.4}{GDP is the same in both. The Jones--Klenow measure, as Kurlat implements it, puts $C+G$ in utility: military spending
is government consumption and counts, while housebuilding is investment and does not. So country A scores higher this year. That
is hard to accept: it assumes a unit of defence is worth a unit of private consumption, and it ignores that the houses will deliver
housing services --- consumption, through rents and imputed rents --- for decades. The measure is a one-year snapshot; investment
pays off to future residents.}
\passo{Jones 9}{(a) Well-being falls: a tenth of the housing stock is lost, and the rebuilding uses labour and materials that would
otherwise have made something new (the broken-window fallacy). (b) Measured GDP probably rises: reconstruction is new residential
investment, and GDP counts the flow of production, not the stock of wealth destroyed. Confirmed: Jones \emph{Instructor's Manual}
(2016), ch.~2, Exercise~8, pp.~12--13.}
\refbib{Kurlat (2020), \S2.2, pp.~33--34 (Jones--Klenow; $C+G$ in utility, p.~34); Jones (2020), ch.~2.}
\rubrica{
\item 1 point per exercise. Any answer that equates GDP with welfare: 0 for that exercise.
\item Kurlat 2.4 answered ``B is better off'' without noting that the measure puts $G$ in and $I$ out: half.
}
\end{sol}

\needspace{12\baselineskip}
% =============================================================================
\bloco{3}{Created: the Prova~01 economy, extended}{Kurlat \S1.1 \,\textbullet\, 35 points}

\questao{(12 points) \emph{Inventories and data sales.} The economy of Prova~01 in a new year. All values in reais;
nothing else happens.
\begin{itemize}[leftmargin=1.4em]
\item Firm A uses no intermediate inputs. It produces inputs worth 800 and pays wages of 500. It sells inputs worth 600 to B and
keeps the remaining 200 in its warehouse.
\item Firm B uses all 600 of the inputs it bought, pays wages of 200, buys advertising from N for 100 and a market-research
dataset from N for 60. It sells final goods to households for 900; final goods worth 150 at market prices are still unsold in its
warehouse at the end of the year.
\item Social network N is free for its users. It sells the advertising (100) and the dataset (60) to B and another dataset to a
firm abroad for 40. It pays wages of 120 and has no other costs.
\end{itemize}}
\noindent\textbf{a)} (7 points) Compute GDP by the product, income and expenditure approaches in one table (by firm for product and
income, by component for expenditure) and check that they agree.\\
\textbf{b)} (3 points) In words: why do the 200 of stored inputs and the 150 of unsold goods count in this year's GDP, and
why does N's data enter the expenditure column once and not the other time?\\
\textbf{c)} (2 points) Next year B sells its 150 of stored goods to households for 150, and nothing else happens. Give the
three approaches for that year and explain the result in one sentence.
\resolespaco
\begin{sol}
\begin{notacao}
\item A, B, N: the input firm, the final-goods firm and the social network. $Q_j$, $Z_j$, $VA_j$, $W_j$, $\Pi_j$ for
$j\in\{A,B,N\}$, and the identities \eqref{q1:G1}--\eqref{q1:G4}: as in Question~1.
\item $C$, $X\ge0$: consumption and exports; $G=M=0$. $I$: investment, here the change in inventories (negative when they are
drawn down).
\end{notacao}

\passo{a) value added}{Gross output is sales plus additions to inventory. A sells 600 and stores 200:}
\begin{align}
Q_A &= (600)+(200) \nonumber\\
&= 800. \label{q7:QA}
\end{align}
A uses no inputs, so by \eqref{q1:G2} and \eqref{q7:QA}:
\begin{align*}
VA_A &= (800)-(0)\\
&= 800 .
\end{align*}
B sells 900 and stores 150:
\begin{align}
Q_B &= (900)+(150) \nonumber\\
&= 1{,}050. \label{q7:QB}
\end{align}
B's intermediate inputs are A's inputs, the advertising and the dataset:
\begin{align}
Z_B &= (600)+(100)+(60) \nonumber\\
&= 760. \label{q7:ZB}
\end{align}
Using \eqref{q1:G2} with \eqref{q7:QB} and \eqref{q7:ZB}:
\begin{align}
VA_B &= (1{,}050)-(760) \nonumber\\
&= 290. \label{q7:vaB}
\end{align}
N sells the advertising and the two datasets, with no inputs:
\begin{align}
Q_N=VA_N &= (100)+(60)+(40) \nonumber\\
&= 200. \label{q7:vaN}
\end{align}
Summing $VA_A=800$, \eqref{q7:vaB} and \eqref{q7:vaN} with \eqref{q1:G1}:
\begin{align*}
\text{GDP} &= (800)+(290)+(200)\\
&= 1{,}290 .
\end{align*}
\passo{a) income}{By \eqref{q1:G3}, each firm's profit is its value added minus its wages:}
\begin{align}
\Pi_A &= (800)-(500) \nonumber\\
&= 300, \label{q7:PiA}\\
\Pi_B &= (1{,}050)-(760)-(200) \nonumber\\
&= 90, \label{q7:PiB}\\
\Pi_N &= (200)-(120) \nonumber\\
&= 80. \label{q7:PiN}
\end{align}
Adding up the wages and the profits \eqref{q7:PiA}--\eqref{q7:PiN}:
\begin{align}
\textstyle\sum_j W_j &= (500)+(200)+(120) \nonumber\\
&= 820, \label{q7:W}\\
\textstyle\sum_j \Pi_j &= (300)+(90)+(80) \nonumber\\
&= 470. \label{q7:Pi}
\end{align}
By \eqref{q1:G3}, the sum of \eqref{q7:W} and \eqref{q7:Pi}:
\begin{align*}
\text{GDP} &= (820)+(470)\\
&= 1{,}290 .
\end{align*}
\passo{a) expenditure}{Investment is the change in inventories, A's 200 plus B's 150:}
\begin{align}
I &= (200)+(150) \nonumber\\
&= 350. \label{q7:I}
\end{align}
With $C=900$, $X=40$ and $G=M=0$, \eqref{q1:G4} and \eqref{q7:I} give
\begin{align*}
\text{GDP} &= (900)+(350)+(40)\\
&= 1{,}290 .
\end{align*}
\begin{center}\footnotesize\begin{tabular}{@{}lr@{\qquad}lr@{\qquad}lr@{}}\toprule
\multicolumn{2}{@{}c}{Production (value added)} & \multicolumn{2}{c}{Income} & \multicolumn{2}{c@{}}{Expenditure}\\\midrule
A: $800-0$ & 800 & A: wages 500 $+$ profit 300 & 800 & Consumption & 900\\
B: $1{,}050-760$ & 290 & B: wages 200 $+$ profit 90 & 290 & Investment (inventories) & 350\\
N: $200-0$ & 200 & N: wages 120 $+$ profit 80 & 200 & Exports & 40\\\midrule
GDP & 1{,}290 & GDP & 1{,}290 & GDP & 1{,}290\\\bottomrule\end{tabular}\end{center}
\passo{b) interpretation}{GDP counts what is \emph{produced} this year, sold or not. Goods produced and kept are an investment in
inventories, valued in the producer's output and in its profit; that is how the product and income columns pick them up. The
dataset sold to B is an input B uses up in production --- intermediate, already inside the price of B's goods --- so it has no
final-use entry. The dataset sold abroad leaves the country: it is a final use by a non-resident, an export. The service users get
for free has no price and stays out, as in Prova~01.}
\passo{c) next year}{Production: B's output is its sales minus the goods drawn from inventory, so by \eqref{q1:G2}:}
\begin{align*}
VA_B=Q_B &= (150)-(150)\\
&= 0 .
\end{align*}
Income: B's revenue 150 against the 150 value of goods leaving the warehouse; by \eqref{q1:G3}:
\begin{align*}
\Pi_B &= (150)-(150)\\
&= 0 .
\end{align*}
Expenditure: $C=+150$ and $I=-150$ (inventories fall); by \eqref{q1:G4}:
\begin{align*}
\text{GDP} &= (150)+(-150)\\
&= 0 .
\end{align*}
GDP $=0$: the goods were counted the year they were made; selling them later moves them from $I$ to $C$ and creates no new
production.
\intuicao{The 160 of N's sales to B and the 600 of A's sales to B appear in the product column only as deductions from B's output;
the final-use column sees only what is left at the end of the chain: consumption, inventories and exports.}
\refbib{Kurlat (2020), \S1.1, Ex.~1.2 (value added), p.~17; Ex.~1.3 (inventories as investment, negative when drawn down), p.~17;
Ex.~1.6 (exports), p.~19. Same treatment of inventories in Williamson (2014), solutions manual, ch.~2, problems 1, 3 and 6
(``additions to inventory are treated as purchasing one's own output'').}
\rubrica{
\item a) 7 points: each approach 2, agreement checked 1. Only one approach: max.\ 2.5 (the Prova~01 Q1(b) error).
\item Inventories left out (GDP $=940$), or the data sold to B added to $C$: max.\ 3 for a).
\item b) 3 points: inventories 1.5, the two data sales 1.5. c) 2 points; GDP of 150 next year: 0 for c).
\item Arithmetic slip with the right layout: $-0.5$, once.
}
\end{sol}

% -----------------------------------------------------------------------------
\questao{(12 points) \emph{Imports, a software sale, dividends and the government.} Another year of the Prova~01 economy.
All values in reais; nothing else happens.
\begin{itemize}[leftmargin=1.4em]
\item Firm A imports raw materials for 100, turns them into inputs and sells the inputs to B for 600. It pays wages of 350.
\item Firm B uses the inputs, buys advertising from N for 100, pays wages of 250 and sells final goods for 1{,}000: 850 to
domestic households and 150 to buyers abroad.
\item Network N is free for its users. It sells the advertising (100) and sells licences of a photo-editing app to domestic
households for 80. It pays wages of 150, of which 100 go to a team that builds a recommendation engine N will use itself for
several years. It has no other costs. Out of its profit, N pays dividends of 40; half of its shareholders live abroad.
\item The government hires content moderators for wages of 30 (their service is free to the public), pays unemployment
benefits of 25 and collects taxes of 55 from households.
\end{itemize}}
\noindent\textbf{a)} (6 points) Compute GDP by the three approaches in one table and check that they agree.\\
\textbf{b)} (2 points) Compute gross national income (GNI) and explain why it differs from GDP.\\
\textbf{c)} (2 points) In words: classify the app licences, the recommendation engine, the dividends and the unemployment
benefits, each in one sentence.\\
\textbf{d)} (2 points) Prova~01 Q1(c) had GDP $=1{,}100$. Account for the difference between that number and your answer
to (a), item by item.
\resolespaco
\begin{sol}
\begin{notacao}
\item A, B, N, and Gov (the government as a producer). $Q_j$, $Z_j$, $VA_j$, $W_j$, $\Pi_j$ (operating surplus for N) and
the identities \eqref{q1:G1}--\eqref{q1:G4}: as in Question~1.
\item $C$, $I$, $G$, $X$, $M\ge0$: consumption, investment, government purchases, exports, imports.
\item GNI: gross national income, the income of residents. $NFI$: net primary income from abroad (negative when more is paid
to non-residents than received from them).
\end{notacao}

\passo{a) value added}{By \eqref{q1:G2}, output minus intermediate inputs. A's input is the imported raw material; B's are A's
inputs and the advertising:}
\begin{align}
VA_A &= (600)-(100) \nonumber\\
&= 500, \label{q8:vaA}\\
VA_B &= (1{,}000)-(600)-(100) \nonumber\\
&= 300. \label{q8:vaB}
\end{align}
N has no inputs; its output is the ads, the app and the own-account software, valued at its cost (the team's wages):
\begin{align}
VA_N &= (100)+(80)+(100) \nonumber\\
&= 280. \label{q8:vaN}
\end{align}
The government's output is valued at cost, $VA_{Gov}=30$. Summing with \eqref{q1:G1}:
\begin{align*}
\text{GDP} &= (500)+(300)+(280)+(30)\\
&= 1{,}110 .
\end{align*}
\passo{a) income}{By \eqref{q1:G3}, profit is value added minus wages, using \eqref{q8:vaA}--\eqref{q8:vaN}:}
\begin{align*}
\Pi_A &= (600)-(100)-(350)\\
&= 150,\\
\Pi_B &= (1{,}000)-(600)-(100)-(250)\\
&= 50,\\
\Pi_N &= (280)-(150)\\
&= 130.
\end{align*}
The government pays wages of 30 and earns no surplus. Adding up wages and profits:
\begin{align}
\textstyle\sum_j W_j &= (350)+(250)+(150)+(30) \nonumber\\
&= 780, \label{q8:W}\\
\textstyle\sum_j \Pi_j &= (150)+(50)+(130) \nonumber\\
&= 330. \label{q8:Pi}
\end{align}
By \eqref{q1:G3}, the sum of \eqref{q8:W} and \eqref{q8:Pi}:
\begin{align}
\text{GDP} &= (780)+(330) \nonumber\\
&= 1{,}110. \label{q8:Y}
\end{align}
\passo{a) expenditure}{Consumption is B's domestic sales plus the app licences:}
\begin{align}
C &= (850)+(80) \nonumber\\
&= 930. \label{q8:C}
\end{align}
With $I=100$ (the recommendation engine), $G=30$, $X=150$ and $M=100$, \eqref{q1:G4} and \eqref{q8:C} give
\begin{align*}
\text{GDP} &= (930)+(100)+(30)+(150)-(100)\\
&= 1{,}110 .
\end{align*}
\begin{center}\footnotesize\begin{tabular}{@{}lr@{\qquad}lr@{\qquad}lr@{}}\toprule
\multicolumn{2}{@{}c}{Production (value added)} & \multicolumn{2}{c}{Income} & \multicolumn{2}{c@{}}{Expenditure}\\\midrule
A: $600-100$ & 500 & A: wages 350 $+$ profit 150 & 500 & Consumption $850+80$ & 930\\
B: $1{,}000-600-100$ & 300 & B: wages 250 $+$ profit 50 & 300 & Investment (software) & 100\\
N: $100+80+100$ & 280 & N: wages 150 $+$ surplus 130 & 280 & Government & 30\\
Government (at cost) & 30 & Government: wages 30 & 30 & Exports & 150\\
 & & & & Imports & $-100$\\\midrule
GDP & 1{,}110 & GDP & 1{,}110 & GDP & 1{,}110\\\bottomrule\end{tabular}\end{center}
\passo{b) GNI}{Half of the dividends is income paid to non-residents out of domestic production, and nothing is received from
abroad:}
\begin{align*}
NFI &= -\frac{(40)}{2}\\
&= -20 .
\end{align*}
GNI adds net primary income from abroad to GDP:
\begin{equation}
\text{GNI}=\text{GDP}+NFI. \label{q8:gni}
\end{equation}
Substitute \eqref{q8:Y} and $NFI=-20$ into \eqref{q8:gni}:
\begin{align*}
\text{GNI} &= (1{,}110)+(-20)\\
&= 1{,}090 .
\end{align*}
GDP measures production inside the country; GNI measures the income of its residents.
\passo{c) classification}{App licences: a service sold to households, final consumption ($C$). Recommendation engine: own-account
software, gross fixed capital formation ($I$), valued at cost. Dividends: a distribution of N's profit already counted in the income
column; they change who receives the income, not GDP (only GNI, for the part paid abroad). Unemployment benefits: a transfer
financed by taxes, no production, in neither $G$ nor GDP; the taxes of 55 do not enter either.}
\passo{d) bridge from 1{,}100}{Start from Prova~01 Q1(c) and add each change, by \eqref{q1:G1}:}
\begin{align*}
\text{GDP} &= (1{,}100)\underbrace{{}+(80)}_{\text{the app, new }C}\ \underbrace{{}+(30)}_{\text{government services}}\
\underbrace{{}-(100)}_{\text{A's imported raw materials}}\\
&= 1{,}110 ,
\end{align*}
which is \eqref{q8:Y}; the 100 of raw materials is no longer domestic value added. Selling 150 of B's goods abroad instead of at
home moves 150 from $C$ to $X$ in \eqref{q1:G4} and leaves GDP unchanged; the dividends, benefits and taxes change nothing.
\intuicao{Imports cancel twice over: the 100 of raw materials is inside the price of B's goods ($C$ and $X$), so it must come out
through $-M$. The software is counted exactly as in Prova~01 Q1(c): as investment, never as the consumption of the workers who
built it --- what they buy with their wages is already in $C$.}
\refbib{Kurlat (2020), \S1.1, Ex.~1.6 (imports and exports), p.~19; Exs.~1.7--1.8 (government at cost, transfers), p.~20; Jones
(2020), ch.~2, p.~22 (intellectual property products in $I$). GDP vs.\ GNI: Williamson (2014), solutions manual, ch.~2, problem~2(c).}
\rubrica{
\item a) 6 points: each approach 1.8, agreement 0.6. Only one approach: max.\ 2.
\item Software counted as consumption: max.\ 3 for a) and 0 for its line in c). Imports not subtracted: max.\ 3 for a).
\item Transfers or taxes added to $G$ or to income: max.\ 3 for a).
\item b) 2 points; GNI $=$ GDP $-40$ (all dividends treated as paid abroad): 1. c) 0.5 per classification. d) 2 points.
}
\end{sol}

% -----------------------------------------------------------------------------
\questao{(11 points) \emph{Capitalise or expense? Depreciation and measured growth.} Two years of the Prova~01 economy. A's
and B's value added is all paid out as factor income (wages and profits; the split is not needed). All values in reais.
\begin{itemize}[leftmargin=1.4em]
\item \textbf{Year 1} is Prova~01 Q1(c): A sells inputs to B for 600; B sells final goods to households for 1{,}000; N sells
advertising to B for 100 and uses that 100 to pay workers who build software for N's own use. The software is ready at the end of
year 1, will be used for four years and depreciates by $100/4=25$ a year from year 2 on.
\item \textbf{Year 2}: A, B and the advertising are as in year 1, and no new software is built. N launches ad-free premium
accounts: households buy accounts for 50, and B buys accounts for 20 for its marketing staff to use at work. N pays its support
staff wages of 30 and has no other costs.
\end{itemize}}
\noindent\textbf{a)} (3 points) Year 1: compute GDP by the three approaches (i) when the software is capitalised (current convention) and
(ii) when it is expensed as a current cost of N (the older convention).\\
\textbf{b)} (3 points) Year 2, software capitalised: compute GDP by the three approaches, and net domestic product (NDP).\\
\textbf{c)} (2 points) Compute GDP growth from year 1 to year 2 under each convention, and explain in words why the sign differs.\\
\textbf{d)} (1.5 points) In words: why do the households' 50 of premium accounts enter $C$ while B's 20 do not?\\
\textbf{e)} (1.5 points) In words: a student wrote that the 100 of software is ``consumption, because the workers spend their
wages''. Explain what is wrong, using the three columns.
\resolespaco
\begin{sol}
\begin{notacao}
\item A, B, N; $Q_j$, $Z_j$, $VA_j$, $W_j$, $\Pi_j$ (operating surplus, gross of depreciation) and the identities
\eqref{q1:G1}--\eqref{q1:G4}: as in Question~1. A superscript $1$ or $2$ marks the year.
\item $C$, $I\ge0$: consumption and investment ($G=X=M=0$). $D_N\ge0$: depreciation of N's software.
\item NDP: net domestic product. $g$: GDP growth rate from year 1 to year 2 (any sign).
\end{notacao}

\passo{a)(i) capitalised}{A's value added is its sales, 600. By \eqref{q1:G2}, B's output minus A's inputs and the advertising:}
\begin{align}
VA_B &= (1{,}000)-(600)-(100) \nonumber\\
&= 300. \label{q9:vaB}
\end{align}
N's output is the ads plus the software, valued at cost:
\begin{align}
VA_N &= (100)+(100) \nonumber\\
&= 200. \label{q9:vaN}
\end{align}
Summing with \eqref{q1:G1}:
\begin{align}
\text{GDP}^1_{\mathrm{cap}} &= (600)+(300)+(200) \nonumber\\
&= 1{,}100. \label{q9:Y1c}
\end{align}
Income: N pays wages of 100, so by \eqref{q1:G3} and \eqref{q9:vaN} its operating surplus is
\begin{align*}
\Pi_N &= (200)-(100)\\
&= 100 .
\end{align*}
Adding A's and B's factor income, N's wages and N's surplus, by \eqref{q1:G3}:
\begin{align*}
\text{GDP} &= (600)+(300)+(100)+(100)\\
&= 1{,}100 .
\end{align*}
Expenditure, by \eqref{q1:G4} with $C=1{,}000$ and $I=100$:
\begin{align*}
\text{GDP} &= (1{,}000)+(100)\\
&= 1{,}100 .
\end{align*}
\passo{a)(ii) expensed}{The software is not output, only a cost, so $VA_N=100$ (the ads). By \eqref{q1:G1}, with $VA_A=600$ and
\eqref{q9:vaB}:}
\begin{align}
\text{GDP}^1_{\mathrm{exp}} &= (600)+(300)+(100) \nonumber\\
&= 1{,}000. \label{q9:Y1e}
\end{align}
Income: N's profit is the ads minus the wages, by \eqref{q1:G3}:
\begin{align*}
\Pi_N &= (100)-(100)\\
&= 0 .
\end{align*}
Adding up incomes, by \eqref{q1:G3}:
\begin{align*}
\text{GDP} &= (600)+(300)+(100)+(0)\\
&= 1{,}000 .
\end{align*}
Expenditure: $C=1{,}000$ and $I=0$, so by \eqref{q1:G4}:
\begin{align*}
\text{GDP} &= (1{,}000)+(0)\\
&= 1{,}000 .
\end{align*}
\begin{center}\footnotesize\begin{tabular}{@{}lrr@{\qquad}lrr@{\qquad}lrr@{}}\toprule
Production & (i) & (ii) & Income & (i) & (ii) & Expenditure & (i) & (ii)\\\midrule
A & 600 & 600 & A & 600 & 600 & Consumption & 1{,}000 & 1{,}000\\
B & 300 & 300 & B & 300 & 300 & Investment & 100 & 0\\
N & 200 & 100 & N: wages & 100 & 100 & & & \\
 & & & N: surplus & 100 & 0 & & & \\\midrule
GDP & 1{,}100 & 1{,}000 & GDP & 1{,}100 & 1{,}000 & GDP & 1{,}100 & 1{,}000\\\bottomrule\end{tabular}\end{center}
\passo{b) year 2}{By \eqref{q1:G2}, B's output minus A's inputs, the advertising and the 20 of accounts:}
\begin{align}
VA^2_B &= (1{,}000)-(600)-(100)-(20) \nonumber\\
&= 280. \label{q9:vaB2}
\end{align}
N has no inputs; its output is the ads plus all premium accounts:
\begin{align}
VA^2_N &= (100)+(50)+(20) \nonumber\\
&= 170. \label{q9:vaN2}
\end{align}
Summing $VA_A=600$, \eqref{q9:vaB2} and \eqref{q9:vaN2} with \eqref{q1:G1}:
\begin{align}
\text{GDP}^2 &= (600)+(280)+(170) \nonumber\\
&= 1{,}050. \label{q9:Y2}
\end{align}
Income: N pays wages of 30, so by \eqref{q1:G3} and \eqref{q9:vaN2} its gross operating surplus is
\begin{align*}
\Pi^2_N &= (170)-(30)\\
&= 140 .
\end{align*}
The software depreciates by a quarter of its cost each year, and the surplus splits into that depreciation and net profit:
\begin{align*}
D_N &= \frac{(100)}{4}\\
&= 25,\\
\Pi^2_N-D_N &= (140)-(25)\\
&= 115 .
\end{align*}
Adding up incomes, by \eqref{q1:G3}:
\begin{align*}
\text{GDP}^2 &= (600)+(280)+(30)+(140)\\
&= 1{,}050 .
\end{align*}
Expenditure: no investment, no exports; by \eqref{q1:G4}:
\begin{align*}
\text{GDP}^2=C &= (1{,}000)+(50)\\
&= 1{,}050 .
\end{align*}
NDP subtracts depreciation from GDP:
\begin{equation}
\text{NDP}=\text{GDP}-D_N. \label{q9:ndp}
\end{equation}
Substitute \eqref{q9:Y2} and $D_N=25$ into \eqref{q9:ndp}:
\begin{align*}
\text{NDP} &= (1{,}050)-(25)\\
&= 1{,}025 .
\end{align*}
\passo{c) growth}{Capitalised: dividing \eqref{q9:Y2} by \eqref{q9:Y1c},}
\begin{align*}
g_{\mathrm{cap}} &= \frac{(1{,}050)}{(1{,}100)}-1\\
&= -4.55\% .
\end{align*}
Expensed: year-2 GDP is also \eqref{q9:Y2} (no software is built in year 2); dividing it by \eqref{q9:Y1e},
\begin{align*}
g_{\mathrm{exp}} &= \frac{(1{,}050)}{(1{,}000)}-1\\
&= +5.00\% .
\end{align*}
Under capitalisation year 1 contains a final good, the software, that year 2 does not produce, so
measured GDP falls although consumption rises from 1{,}000 to 1{,}050; under expensing the software never appears, and GDP only
sees the rise in consumption. The convention changes the level of GDP in the year the asset is built and, through it, the growth rate.
\passo{d) intermediate vs final}{The households use the accounts up themselves: final consumption, $C$ in \eqref{q1:G4}. B uses the
same accounts as an input for producing and selling its goods: intermediate consumption, in $Z_B$ of \eqref{q1:G2}, so B's value
added falls from \eqref{q9:vaB} to \eqref{q9:vaB2}:}
\begin{align*}
VA_B-VA^2_B &= (300)-(280)\\
&= 20 .
\end{align*}
That 20 is already inside the price of B's goods; adding it to $C$ would count it twice. N's value added rises by the full 70 in
both cases ($50+20$ in \eqref{q9:vaN2}).
\passo{e) the exam error}{The wages are on the \emph{income} side, \eqref{q1:G3}: N's workers earn 100, and what they buy with it
is part of $C=1{,}000$, which the statement keeps unchanged. The expenditure side, \eqref{q1:G4}, asks who uses the \emph{output}:
the software is kept by N and used for four years, so it is investment. Calling it consumption counts the workers' income as a final use and, with $C$ fixed at 1{,}000, leaves
the expenditure column at 1{,}000 against 1{,}100 in the other two --- the three columns stop agreeing, which flags the error.}
\intuicao{Capitalising treats own-account software like a machine bought from outside: an asset produced now that pays off later, its
wear recorded as depreciation (in NDP, not in GDP). Expensing treats it as a cost, like electricity. Same economy, different GDP.}
\refbib{Kurlat (2020), \S1.1, Ex.~1.3 (investment), p.~17; Ex.~1.10 (depreciation; GDP is gross), p.~21. Jones (2020), ch.~2,
p.~22 (intellectual property products --- software, R\&D --- in $I$) and p.~25 (effect on the labor share).}
\rubrica{
\item a) 3 points: 1.5 per convention, the three approaches each. Only GDP totals: max.\ 1.
\item b) 3 points: three approaches 2.4, NDP 0.6. Depreciation subtracted from GDP: max.\ 1.5 for b).
\item c) 2 points: both rates 1, explanation 1. d) 1.5 points. e) 1.5 points; ``it is consumption'' repeated: 0 for e).
}
\end{sol}


\end{document}
